% appC1 —— 附录 C 前段：积分方程的渐近（有限维情形）（原文 wcf-paper.tex 3529–3685 行；共 10 段、13 条编号公式 (C.1)–(C.13)）
\section{积分方程的渐近}
\label{app:integral-asymptotics}

本附录中，我们先说明如何修改 \cite{Cecotti:1993rm} 的渐近分析，使之适用于存在多条 BPS 射线的情形；然后把这个分析改造成适合我们问题的形式。

\subsection*{有限维情形}

文献 \cite{Cecotti:1993rm} 研究的是复 $x$ 平面上的一个黎曼--希尔伯特问题：关于 $m \times m$ 矩阵 $\Psi(x)$，其不连续性沿实轴分布：
\begin{align}
\Psi(y e^{- i \epsilon}) &= \Psi(y e^{ i \epsilon}) S \qquad \text{当 } y \in
\IR^+, \notag \\
\Psi(y e^{- i \epsilon}) &= \Psi(y e^{ i \epsilon}) S^t \qquad \text{当 } y \in
\IR^-. \label{eq:cv-discont}
\end{align}
（此处我们已回到算符复合的标准约定。）这里充当中心荷 $Z_\gamma$ 之对应物的是复数 $\Delta_{ij}$（$i,j = 1, \dots, m$），它们满足 $\Delta_{ij} = w_i-w_j$，其中诸 $w_i$ 为某些复数。矩阵 $S$ 是三角矩阵，对角元为 $1$，并且当 $\re \Delta_{ij} <0$ 时 $S_{ij}=0$。

$\Psi$ 的渐近行为由诸常数 $w_i$ 决定。若定义
\begin{equation} \label{eq:appC-Phi-def}
\Phi_{ij}(x) = \Psi_{ij}(x) e^{- \beta x w_j - \frac{\beta}{x} \bar w_j},
\end{equation}
则 $\Phi(x)$ 在 $x\to \infty$ 时趋于单位矩阵，在 $x\to 0$ 时趋于某个「度量」$g_{ij}$。

矩阵 $S$ 是该问题的「斯托克斯乘子」。把它重新表示成若干个更基本的「斯托克斯因子」之乘积，会比较方便。事实上，只要假定没有三个 $w_i$ 在复平面中共线，就存在唯一的分解
\begin{equation}\label{eq:appC-S-decomp}
S = \prod^\cwarrow_{(ij): \re \Delta_{ij} > 0} s_{(ij)}, \quad S^t = \prod^\cwarrow_{(ij): \re \Delta_{ij} < 0} s_{(ji)}
\end{equation}
其中两个乘积都按 $\arg \Delta_{ij}$ 递增的次序取，而每个 $s_{(ij)}$ 的对角元为 $1$，并且在位置 $(ji)$ 处有唯一的非零非对角元，其值为 $- \mu_{ij} = -\mu_{ji}$。

利用这一分解，我们可以引入黎曼--希尔伯特问题的「多射线」版本。具体地说，在 $x$ 平面上引入一组过原点的射线
\begin{equation}\label{eq:appC-rays}
\ell_{(ij)} = \{ x: x \Delta_{ij} \in \IR_+ \},
\end{equation}
并要求
\begin{equation}\label{eq:appC-jump}
\Psi(y e^{- i \epsilon}) = \Psi(y e^{ i \epsilon}) s_{(ij)} \qquad \text{当 } y \in
\ell_{(ij)}.
\end{equation}
重要的是，可以证明：无论采用单射线问题还是多射线问题，$g_{ij}$——它在文献 \cite{Cecotti:1993rm} 中正是主要的研究对象——都相同。

文献 \cite{Cecotti:1993rm} 中关于 $\Phi(x)$ 的积分方程 (4.17) 为（其中对 $\epsilon$ 的约定作了少许修改，以便后文使用）：
\begin{align}
\Phi_{ij}(x) = \delta_{ij} &+ \frac{1}{2 \pi i} \int_0^\infty
\frac{dy}{y-x} \sum_k \Phi_{ik}(e^{i \epsilon} y)(1-S)_{kj}e^{-
\beta y \Delta_{kj} - \beta/y \bar \Delta_{kj}}  \nonumber  \\
&+ \frac{1}{2 \pi i}
\int_{-\infty}^0 \frac{dy}{y-x} \sum_k \Phi_{ik}(e^{-i \epsilon}
y)(1-S^t)_{kj}e^{- \beta y \Delta_{kj} - \beta/y \bar \Delta_{kj}}. \label{eq:cv-integral}
\end{align}
该方程的一个解给出单射线黎曼--希尔伯特问题的一个解。现在我们表述一个新的积分方程，它与多射线问题等价：
\begin{equation} \label{eq:cv-multi-line}
\Phi(x)_{ij} = \delta_{ij} + \sum_{k} \frac{1}{2 \pi i}
\int_{\ell_{(kj)}} \frac{dy}{y-x}\Phi_{ik}(y)\mu_{kj}e^{- \beta y
\Delta_{kj} - \frac{\beta}{y} \bar \Delta_{kj}}
\end{equation}
这里的积分应理解为主值积分。注意：当没有三个 $w_i$ 共线时，矩阵元 $\Phi_{ik}$ 在穿过所有射线 $\ell_{(kj)}$ 时都是连续的。

\eqref{eq:cv-multi-line} 的递推解取一个简单的形式：
\begin{equation}\label{eq:appC-Phi-series}
\Phi(x)_{ij} = \delta_{ij} + \Phi(x)_{ij}^{(1)} + \Phi(x)_{ij}^{(2)} + \cdots,
\end{equation}
其中
\begin{align}
\Phi(x)_{ij}^{(1)} &= \frac{1}{2 \pi i} \mu_{ij}
\int_{\ell_{(ij)}} \frac{dy}{y-x} e^{ -\beta y
\Delta_{ij} - \frac{\beta}{y} \bar \Delta_{ij} }, \label{eq:appC-Phi12}\\
\Phi(x)_{ij}^{(2)} &= \frac{1}{(2 \pi i)^2} \sum_{i_2} \mu_{i i_2} \mu_{i_2 j}
\int_{\ell_{(i i_2)}} \int_{\ell_{(i_2 j)}} \frac{dy_1}{y_1-y_2}
\frac{dy_2}{y_2-x}e^{ -\beta
y_1 \Delta_{i i_2} - \frac{\beta}{y_1} \bar \Delta_{i i_2} }e^{ -\beta
y_2 \Delta_{i_2 j} - \frac{\beta}{y_2} \bar \Delta_{i_2 j} }, \notag
\end{align}
一般地，$\Phi(x)^{(n)}$ 涉及沿 $n$ 条射线 $\ell_{(i_k i_{k+1})}$ 的所有链的积分，其中 $i_1 = i$、$i_{n+1} = j$。
沿 $\ell_{(i_k i_{k+1})}$ 的每个积分都含有因子 $e^{ - \beta y \Delta_{i_k i_{k+1}} - \frac{\beta}{y} \bar \Delta_{i_k i_{k+1}} }$。在这条射线上，该指数是实的且为负，并在 $y = \exp (-i \arg \Delta_{i_k i_{k+1}})$ 处有唯一的峰。当 $\beta$ 取得很大时，该积分便由鞍点法很好地近似：把被积函数的其余部分换成它在峰点处的值。于是，固定 $x$ 时的大 $\beta$ 渐近式就简单地是：\footnote{这些渐近式成立，除非 $x$ 恰好落在对 $y_n$ 积分的鞍点上，即 $x = \exp [-i \arg \Delta_{i_n j}]$。在这一点上，我们得到类似的大 $\beta$ 渐近式，只是诸 $\sqrt{\beta}$ 压低因子中少了一个。}
\begin{equation}\label{eq:appC-asym1}
\Phi(x)^{(1)}_{ij}  \sim \frac{1}{2 \sqrt{\pi \beta |\Delta_{ij}|} i}
\mu_{ij}\frac{1}{\exp[-i \arg \Delta_{ij}]-x} e^{- 2 \beta
|\Delta_{ij}|},
\end{equation}
以及
\begin{equation}\label{eq:appC-asym2}
\begin{aligned}
\Phi(x)^{(2)}_{ij} &\sim \sum_{i_2} \mu_{i i_2} \mu_{i_2 j} \frac{1}{2 \sqrt{\pi \beta
|\Delta_{i i_2}|} i}\frac{1}{2 \sqrt{\pi \beta |\Delta_{i_2 j}|} i} \times \\
&\qquad \frac{1}{\exp[-i \arg \Delta_{i_2 j}]-x}\frac{1}{\exp[-i \arg
\Delta_{i i_2}]-\exp[-i \arg \Delta_{i_2 j}]}e^{- 2 \beta |\Delta_{i
i_2}|}e^{- 2 \beta |\Delta_{i_2 j}|},
\end{aligned}
\end{equation}
更高阶的 $\Phi^{(n)}$ 有类似的估计。由三角不等式可见，在大 $\beta$ 极限下 $\Phi^{(2)}$ 相对于 $\Phi^{(1)}$ 是指数压低的；同样，$\Phi^{(n+1)}$ 相对于 $\Phi^{(n)}$ 也是压低的。$\Phi^{(n)}$ 所具有的指数压低，恰好就是 $n$-瞬子修正所应期待的量级。

我们注意到，这一渐近分析比 \eqref{eq:cv-integral} 的相应分析要简单得多；在那种情形，必须变形积分围道使其穿过相应的鞍点，而途中会遇到割线与极点，其贡献必须仔细追踪。

为结束本节，我们简要讨论这一展开的解析性质。如上所见，$\Phi$ 的第 $n$ 阶修正可以用某些迭代积分表示：
\begin{equation}\label{eq:appC-Fn}
\begin{aligned}
\CF^{(n)}[x,\beta;\Delta_{i_1 i_2},\dots, \Delta_{i_{n} i_{n+1}}] ={}& \\
&\prod_{k=1}^n \left[ \int_{\ell_{(i_k i_{k+1})}} \frac{dy_k}{2 \pi i} e^{ -\beta y_k \Delta_{i_k i_{k+1}} -
\frac{\beta}{y_k} \bar \Delta_{i_k i_{k+1}} } \right] \\
&\times \prod_{k=1}^n
\frac{1}{y_k - y_{k+1}} \Bigg\vert_{y_{n+1} = x}.
\end{aligned}
\end{equation}
这个 $\CF^{(n)}$ 在射线 $\arg x = -\arg \Delta_{i_n i_{n+1}}$ 上有明显的不连续性，该不连续性等于
\begin{equation}\label{eq:appC-disc}
\CF^{(n-1)}[x,\beta;\Delta_{i_1 i_2}, \dots, \Delta_{i_{n-1}
i_{n}}]e^{ - \beta x \Delta_{i_{n} i_{n+1}} - \frac{\beta}{x} \bar
\Delta_{i_{n} i_{n+1}} }.
\end{equation}
注意，$\CF^{(n)}$ 对自变量 $d_k = \Delta_{i_k i_{k+1}}$ 的一般取值是有意义的，不必受 $\Delta_{ij} = w_i-w_j$ 这一限制。（事实上，甚至可以把 $\Delta_{ij}$ 与 $\bar \Delta_{ij}$ 取作彼此独立的复参数，而把诸射线放在 $x^2 \Delta_{ij}/\bar \Delta_{ij} \in \IR^+$ 处。）只要两个相邻自变量 $d_k, d_{k+1}$ 的相位对齐，它就有割线。此时的不连续性就是取同样自变量的 $\CF^{(n-1)}$，只须作替换 $d_k, d_{k+1} \to d_k + d_{k+1}$。

在文献 \cite{Cecotti:1993rm} 的渐近分析中，出现过本质上相同的函数 $\CF^{(n)}$；我们的函数与它们只差在分支割线的位置。
